\documentclass[10pt,a4paper]{amsart}
\usepackage[margin=19mm]{geometry}
\usepackage{amsmath,amssymb,mathtools}
\usepackage[scr=rsfs,cal=euler]{mathalfa}
\usepackage{xfrac}
\usepackage[unicode,hidelinks,bookmarks=false]{hyperref}
\newcommand{\ie}{i.e.\ }
\newcommand{\cf}{cf.\ }
\newcommand{\resp}{respectively\ }
\newcommand{\Subsection}{\subsection}
\providecommand{\nobreakdash}{\nobreak\hbox{-}}
\setlength{\emergencystretch}{2em}
\multlinegap0pt
\usepackage{xcolor}
\usepackage{amssymb}
\usepackage{tikz}
\newtheorem{definition}{Definition}
\newtheorem{proposition}{Proposition}
\newtheorem{theorem}{Theorem}
\title{Spectral separation of variables from equivalent Lagrangian systems}
\author{Mattia Scomparin}
\date{Source: arXiv:2605.15679v1 (CC BY 4.0)}
\begin{document}
\maketitle

\noindent\emph{Source and licence.} Spectral separation of variables from equivalent Lagrangian systems. Version arXiv:2605.15679v1 (CC BY 4.0), distributed by arXiv under the Creative Commons Attribution 4.0 International licence, http://creativecommons.org/licenses/by/4.0/, as recorded in the arXiv licence metadata for this exact version and shown on the abstract page https://arxiv.org/abs/2605.15679v1. Full licence notice: \texttt{licenses/arxiv-2605.15679-CC-BY-4.0.txt}.\par\smallskip

\noindent\textit{Editorial scope.} Counted unit: the article's theorem th:weak (source0.tex lines 623--648). The article's complete original proof of it is delivered with it (source0.tex lines 650--694, one \texttt{proof} environment of 45 lines). No mathematics is written, repaired or re-derived: the statement and the proof are the article's own lines, in the article's own notation, and no retained line is substituted or re-flowed.  Retained uncounted as the source-local context the proof needs: lines 419--425; lines 465--470; lines 567--586; lines 591--612.  Nothing was omitted from any retained span, so the package carries no omission record.  No retained line contains a citation, so the wrapper carries no bibliography and no bibliography entry.  No retained line contains a figure environment, an image-inclusion command or drawing code, so no image is delivered and the package declares no asset dependency.  Editorial formatting only, in addition: an AMS \texttt{amsart} preamble that restates the article's own declarations and macros used by the retained lines, namely the environments \texttt{definition}, \texttt{proposition}, \texttt{theorem} on separate counters, together with the standard packages those lines need.  The retained environments print in this document as Definition 1, Proposition 1, Theorem 1 in order of appearance, whereas the article prints the same items with the numbering of its own full document.  The complete native source of the article as distributed in the arXiv e-print for this exact version is bundled unchanged as \texttt{sources/200671/source0.tex}.
\begin{equation}
\label{eq:xEoM}
\ddot{x}_k + \partial_{x_k}U(x) = 0,
\qquad
\lambda_k \ddot{x}_k + \partial_{x_k} \widetilde U(x) = 0,
\qquad k = 1,\dots,n.
\end{equation}

\begin{equation}
\label{eq:energies}
E(x,\dot{x})=\frac{1}{2} \sum_{k=1}^n \dot{x}_k^2+U(x),
\quad\quad
E(x,\dot{x})=\frac{1}{2} \sum_{k=1}^n \lambda_k \dot{x}_k^2+\widetilde U(x)\,.
\end{equation}

\begin{definition}[Spectral framework]
\label{def:spectral}
Let $A \in \mathrm{Sym}(n,\mathbb{R})$. We adopt the following spectral framework:
\begin{enumerate}
\item Let $Q \in O(n)$ diagonalize $A$, that is
$
Q^T A Q = \mathrm{diag}(\lambda_1 I_{m_1}, \dots, \lambda_r I_{m_r}),
$
where $\lambda_1,\dots,\lambda_r$ are the distinct eigenvalues of $A$ with
multiplicities $m_1,\dots,m_r$.
\item Decompose $Q = [Q_{(1)}, \dots, Q_{(r)}]$, where the columns of
$Q_{(k)} \in \mathbb{R}^{n \times m_k}$ form an orthonormal basis of the eigenspace
$\mathcal{E}_{\lambda_k}$.
\item Define the spectral coordinates $x := Q^T q \in \mathbb{R}^n$, with
$
x_{(k)} := Q_{(k)}^T q \in \mathbb{R}^{m_k},
$
so that $x = (x_{(1)}, \dots, x_{(r)})^T$.
\end{enumerate}
\end{definition}

\begin{proposition}[Equivalence conditions]
\label{theo:eq_statmnt}
Assume the spectral framework of Definition~\ref{def:spectral}. The following
statements are equivalent:
\begin{enumerate}
\item
The Lagrangians
$
\mathcal{L}(q,\dot q)=\tfrac12|\dot q|^2-U(q)
$
and
$
\widetilde{\mathcal{L}}(q,\dot q)=\tfrac12 \dot q^T A \dot q-\widetilde U(q)
$
generate the same Euler--Lagrange equations.
\item
The potentials satisfy
$
\partial_{x_i}\widetilde U=\lambda_i\,\partial_{x_i}U$, with $i=1,\dots,n.
$
\end{enumerate}
\end{proposition}

\begin{theorem}[Weak spectral separation and first integrals]
\label{th:weak}
Under the assumptions of Proposition~\ref{theo:eq_statmnt}, the potentials split
along the spectral decomposition of $A$ as
\begin{equation}
\label{eq:weak-potentials}
U=\sum_{k=1}^r H_k[x_{(k)}],
\qquad
\widetilde U=\sum_{k=1}^r \lambda_k H_k[x_{(k)}]+C,
\end{equation}
for suitable smooth functions $H_k:\mathbb{R}^{m_k}\to\mathbb{R}$. Consequently,
the Euler--Lagrange equations decouple into $r$ independent blocks,
\begin{equation}
\label{eq:weak-eom}
\ddot x_{(k)}+\nabla H_k[x_{(k)}]=0,
\qquad k=1,\dots,r.
\end{equation}
Moreover, the corresponding block energy sums
\begin{equation}
\label{eq:weak-energies}
E=\sum_{k=1}^r\!\Big(\tfrac12|\dot x_{(k)}|^2+H_k[x_{(k)}]\Big),
\quad
\widetilde E=\sum_{k=1}^r \lambda_k\!\Big(\tfrac12|\dot x_{(k)}|^2+H_k[x_{(k)}]\Big)+C,
\end{equation}
are conserved along solutions.
\end{theorem}

\begin{proof}
%
$[$Step $1]$
Let $x_\alpha \in (\lambda_a, \mathcal{E}_{\lambda_a})$ and $x_\beta \in (\lambda_b, \mathcal{E}_{\lambda_b})$ be spectral coordinates associated with two distinct block eigenspaces ($\lambda_a \neq \lambda_b$). Using assumption (2) of Proposition \ref{theo:eq_statmnt}, differentiate $\partial_{x_\alpha} \widetilde U = \lambda_a \partial_{x_\alpha} U$ with respect to $x_\beta$ to obtain
$
\partial^2_{x_\beta x_\alpha} \widetilde U = \lambda_a \partial^2_{x_\beta x_\alpha} U.
$
Similarly, differentiating $\partial_{x_\beta} \widetilde U = \lambda_b \partial_{x_\beta} U$ with respect to $x_\alpha$ yields
$
\partial^2_{x_\alpha x_\beta} \widetilde U = \lambda_b \partial^2_{x_\alpha x_\beta} U.
$
Since $U, \widetilde U \in C^2(\mathbb{R}^n)$, mixed partial derivatives commute, hence
$
\partial^2_{x_\alpha x_\beta} \widetilde U = \partial^2_{x_\beta x_\alpha} \widetilde U.
$
Therefore,
$
\lambda_a \partial^2_{x_\alpha x_\beta} U = \lambda_b \partial^2_{x_\alpha x_\beta} U,
$
which implies
$
(\lambda_a - \lambda_b)\, \partial^2_{x_\alpha x_\beta} U = 0.
$
Since $\lambda_a \neq \lambda_b$, it follows that
$
\partial^2_{x_\alpha x_\beta} U = 0.
$
By the same argument, one also obtains $\partial^2_{x_\alpha x_\beta} \widetilde U = 0$. Thus, all mixed derivatives between coordinates belonging to distinct eigenspaces vanish.
$[$Step $2]$
Since $\partial^2_{x_\alpha x_\beta} U = 0$ and $\partial^2_{x_\alpha x_\beta} \widetilde U = 0$, integration of these conditions implies the existence of smooth functions
$
H_k : \mathbb{R}^{m_k} \to \mathbb{R}$ and 
$\widetilde H_k : \mathbb{R}^{m_k} \to \mathbb{R},
$
such that
$
U(x) = H_1[x_{(1)}] + \cdots + H_r[x_{(r)}]$ and 
$\widetilde U(x) = \widetilde H_1[x_{(1)}] + \cdots + \widetilde H_r[x_{(r)}].
$
Using assumption (2) of Proposition \ref{theo:eq_statmnt}, namely $\partial_{x_\alpha} \widetilde U = \lambda_a \partial_{x_\alpha} U$, we deduce that
$
\widetilde H_k[x_{(k)}] = \lambda_k\, H_k[x_{(k)}].
$
Substituting these expressions into Eqs.~\eqref{eq:xEoM} and \eqref{eq:energies}, the desired results follow.
\end{proof}

\end{document}
